% Generated from paper/duality_confinement/main.tex by scripts/sync_paper.py.
% Edit the modular sources and regenerate this copy.
\documentclass[11pt]{article}

\usepackage[T1]{fontenc}
\usepackage[utf8]{inputenc}
\usepackage[english]{babel}
\usepackage[margin=1in]{geometry}
\usepackage{lmodern}
\usepackage{microtype}
\usepackage{mathtools}
\usepackage{amsmath,amssymb,amsthm}
\usepackage{array}
\usepackage{booktabs}
\usepackage{tabularx}
\usepackage{longtable}
\usepackage{needspace}
\usepackage{float}
\usepackage{graphicx}
\usepackage{pdflscape}
\usepackage{caption}
\usepackage{enumitem}
\usepackage{fancyhdr}
\usepackage[numbers]{natbib}
\usepackage{doi}
\usepackage{xurl}
\usepackage{aliascnt}
\usepackage{hyperref}
\usepackage{cleveref}
\usepackage{orcidlink}
\usepackage{titling}



\DeclareMathOperator{\Fix}{Fix}
\DeclareMathOperator{\Ran}{Ran}
\DeclareMathOperator{\dist}{dist}
\DeclareMathOperator{\trace}{tr}

\newcommand{\Pminus}{P_{-}}
\newcommand{\psim}{\psi_{-}}
\newcommand{\AX}{A_{X}}
\newcommand{\Kminus}{K^{-}}
\newcommand{\Jiso}{J_{\mathrm{iso}}}
\newcommand{\IOL}{\hat{\mathfrak{X}}}

\newcommand{\GamSDTC}{\Gamma_{\mathrm{SDTC}}}

\newcommand{\SiblingRH}{[2]}

\newtheorem{theorem}{Theorem}[section]
\newcommand{\sharedtheorem}[2]{\newaliascnt{#1}{theorem}\newtheorem{#1}[#1]{#2}\aliascntresetthe{#1}}
\sharedtheorem{lemma}{Lemma}
\sharedtheorem{proposition}{Proposition}
\sharedtheorem{corollary}{Corollary}
\theoremstyle{definition}
\sharedtheorem{definition}{Definition}
\sharedtheorem{remark}{Remark}
\newtheorem*{example}{Example}
\floatstyle{ruled}
\newfloat{algorithm}{tbp}{loa}
\floatname{algorithm}{Algorithm}

\urlstyle{same}
\makeatletter
\g@addto@macro\UrlBreaks{\do\-\do\:}
\makeatother
\captionsetup{
  font=small,
  labelfont=bf,
  labelsep=period
}
\captionsetup[table]{skip=6pt,position=top}
\captionsetup[figure]{skip=6pt}
\captionsetup[algorithm]{skip=6pt,position=top}
\crefname{algorithm}{algorithm}{algorithms}
\Crefname{algorithm}{Algorithm}{Algorithms}
\crefname{theorem}{theorem}{theorems}
\Crefname{theorem}{Theorem}{Theorems}
\crefname{lemma}{lemma}{lemmas}
\Crefname{lemma}{Lemma}{Lemmas}
\crefname{proposition}{proposition}{propositions}
\Crefname{proposition}{Proposition}{Propositions}
\crefname{corollary}{corollary}{corollaries}
\Crefname{corollary}{Corollary}{Corollaries}
\crefname{definition}{definition}{definitions}
\Crefname{definition}{Definition}{Definitions}
\crefname{remark}{remark}{remarks}
\crefname{appendix}{Appendix}{Appendices}
\Crefname{appendix}{Appendix}{Appendices}
\Crefname{remark}{Remark}{Remarks}
\setlist[enumerate]{leftmargin=*, itemsep=2pt, topsep=4pt}
\setlength{\droptitle}{-3.2em}
\pretitle{\begin{center}\LARGE\bfseries}
\posttitle{\par\end{center}\vspace{0.5em}}
\preauthor{\begin{center}\normalsize}
\postauthor{\par\end{center}\vspace{-0.4em}}
\predate{\begin{center}\small}
\postdate{\par\end{center}\vspace{0.6em}}
\setlength{\emergencystretch}{2em}
\clubpenalty=10000
\widowpenalty=10000
\displaywidowpenalty=10000
\interfootnotelinepenalty=10000
\allowdisplaybreaks[2]
\fancypagestyle{firstpage}{\fancyhf{}
  \fancyfoot[C]{\scriptsize Automorph Inc., Wilmington, DE, USA \quad \texttt{ioannis@automorph.io} \quad DOI: 10.5281/zenodo.23086972\\[1pt]\textcopyright\ 2026 \quad CC-BY 4.0 \quad Preprint --- not peer reviewed.}
  \renewcommand{\headrulewidth}{0pt}
  \renewcommand{\footrulewidth}{0pt}
}

\title{Self-Dual Trace Confinement: A Six Birds Structural Law for Formed Closures Under Involutive Self-Duality}
\author{Ioannis Tsiokos\,\orcidlink{0009-0009-7659-5964}}
\date{Version 3: October 1, 2026\\[2pt] \footnotesize (Version 1: June 17, 2026)}
\hypersetup{
  pdftitle={Self-Dual Trace Confinement: A Six Birds Structural Law for Formed Closures Under Involutive Self-Duality},
  pdfauthor={Ioannis Tsiokos},
  hidelinks
}

\DeclareRobustCommand{\leanid}[1]{{\def\UrlBreaks{\do\.\do\_}\def\UrlBigBreaks{}\def\UrlNoBreaks{}\nolinkurl{#1}}}
\begin{document}

\maketitle
\thispagestyle{firstpage}

\begin{abstract}
Let \(J\) be an involution of a measure space \((X,\mu)\), and let
\(\psi\) be a strongly measurable observable with values in a real
Hilbert space \(Y\) that intertwines \(J\) with a linear isometric
involution of \(Y\).  The odd part \(\psim\) of \(\psi\) vanishes on
the fixed locus \(\Fix(J)\); it is called separating when it vanishes
only there.  When \(\int_X\|\psim\|^2\,d\mu<\infty\), the Gram operator
\(\AX=\int_X\psim\psim^*\,d\mu\) is positive and trace-class with
\(\trace\AX=\int_X\|\psim\|^2\,d\mu\).  Consequently \(\AX\) is
dominated in the Loewner order by positive trace-class operators
\(B_n\) with \(\trace B_n\to0\) if and only if \(\psim=0\) almost
everywhere, and then, for a separating readout, \(\mu\)-almost every
point is fixed by \(J\).  The paper records this equivalence together
with a Markov-type bound under quantitative separation, variants for
domination on finite windows and for budgets with two defect terms,
and Douglas's characterization of exact domination.  The mathematics is
elementary.  Its purpose is to fix the precise form of a hypothesis in
the author's Six Birds framework, Self-Dual Trace Confinement (SDTC):
that in specified self-dual settings such domination records are
supplied by independent structure, for instance by an operator built
from other data.  SDTC is a named hypothesis, not a theorem: for a
given ledger with a separating odd readout it is equivalent to
confinement, and it cannot hold for all involutive ledgers, as a
complex-conjugation example shows.  All theorems, lemmas and
propositions are formalized in Lean~4 with Mathlib for real Hilbert
spaces of arbitrary dimension, with the trace of a positive operator
taken as the sum of its diagonal entries in a Hilbert basis and the
Gram operator as a Bochner integral in the operator norm; the
trace-class integrability of the integrand and the standard facts
identifying this trace with the trace-class trace are used only in the
text.
\end{abstract}

\medskip
\noindent\textbf{Keywords.} involution; fixed locus; positive operators; trace-class operators; Loewner order; Douglas factorization; Self-Dual Trace Confinement; Six Birds Theory; Lean formalization.

\section{Introduction}
\label{sec:intro}

When must an obstruction attached to a symmetry vanish?  Here is the
form of the question studied in this paper.  A space of objects \(X\)
carries an involution \(J\), so that applying \(J\) twice returns the
original object, and a measure \(\mu\) that records which objects are
present.  The objects are observed through a strongly measurable map
\(\psi\) with values in a real Hilbert space \(Y\), and \(Y\) carries a
linear isometric involution \(\Jiso\) compatible with \(\psi\) in the
sense that \(\psi(Jx)=\Jiso\psi(x)\).  The involution \(\Jiso\) splits every
response into an even part and an odd part.  The odd part is given by
the orthogonal projection
\[
  \Pminus=\tfrac12\,(I_Y-\Jiso),
  \qquad
  \psim(x)=\Pminus\psi(x),
\]
and it changes sign under the duality: \(\psim(Jx)=-\psim(x)\).  In
particular \(\psim\) vanishes at every fixed point of \(J\).  When it
vanishes \emph{only} there, the readout \(\psim\) detects exactly
the displacement of an object from the fixed locus \(\Fix(J)\).

The obstruction is the positive operator
\[
  \AX=\int_X \psim(x)\,\psim(x)^*\,d\mu(x),
\]
where \(v^*\) denotes the functional \(h\mapsto\langle v,h\rangle\),
so that \(vv^*\) is the rank-one operator
\(h\mapsto\langle v,h\rangle v\).  When
\(\int_X\|\psim\|^2\,d\mu<\infty\), the operator \(\AX\) is positive
and trace-class, and
\[
  \trace\AX=\int_X\|\psim(x)\|^2\,d\mu(x).
\]
Suppose now that \(\AX\) is dominated by positive trace-class
operators whose traces tend to zero:
\[
  \AX\preceq B_n\quad(n\geq1),
  \qquad
  \trace B_n\to0,
\]
where \(\preceq\) is the Loewner order (\(A\preceq B\) means that
\(B-A\) is positive).  Monotonicity of the trace gives
\(0\leq\trace\AX\leq\trace B_n\) for every \(n\), so
\(\trace\AX=0\).  Hence \(\AX=0\) and \(\psim=0\) almost everywhere.
If \(\psim\) vanishes only on \(\Fix(J)\), then \(\mu\)-almost every
object is fixed by \(J\): the visible mass is \emph{confined} to the
fixed locus.  Conversely, if \(\psim=0\) almost everywhere then
\(\AX=0\) and \(B_n=0\) will do.  The equivalence is immediate; it is
recorded as \Cref{thm:duality_confinement:master-theorem}.

The argument is short, and none of its operator-theoretic ingredients
is new: the Loewner order \citep{Loewner1934,Bhatia1997}, trace
ideals \citep{Simon2005}, Markov's inequality \citep{Folland1999},
the arithmetic--geometric mean inequality
\citep{HardyLittlewoodPolya1952} and Douglas's factorization theorem
\citep{Douglas1966} are classical.  Nor does the theorem make
confinement easier to prove: by the equivalence just stated, producing
domination records is exactly as hard as proving
\(\int_X\|\psim\|^2\,d\mu=0\).  What the formulation offers is a
format for the case in which \(\AX\) is not accessible directly.  In
applications of the framework one computes, from other data, a
positive operator \(\Kminus\) on \(Y\)---for instance a carrier-side
quantity built, in finite dimensions, from an audit operator and a
probe \citep{TsiokosNeedleKiller2026}---and establishes, by a separate
argument, a bridge inequality \(\AX\preceq\Kminus+E\) with an explicit
positive defect \(E\).  Trace bounds on \(\Kminus\) and \(E\) then
control \(\trace\AX=\int_X\|\psim\|^2\,d\mu\), the odd energy of the
ledger; under quantitative separation this in turn bounds the mass at
a given distance from the fixed locus.
All of the substance lies in the bridge; the results of this paper
record what follows from it.

The paper therefore separates two things that are easy to conflate:
\begin{enumerate}
  \item the \emph{confinement theorem} and its variants, proved here,
    which convert domination records into confinement on the fixed
    locus; and
  \item a \emph{hypothesis} about which structures possess such
    records.  In the author's Six Birds framework
    \citep{TsiokosFoundationsI2026} this hypothesis is named
    Self-Dual Trace Confinement, written \(\GamSDTC\) (SDTC for
    short): in specified closures that carry a genuine involutive
    self-duality, trace-vanishing domination records for the odd
    readout are supplied by the structure.  SDTC is not proved here,
    and no result of this paper depends on it.  It is not a
    statement about all involutive ledgers: the first example in
    \Cref{sec:involutive_ledger}, taken with Lebesgue measure on a
    disc centred on the real axis, is an involutive ledger with a
    separating odd readout that is not confined, and hence has no
    such records.
\end{enumerate}
\Cref{sec:framework} explains the framework vocabulary in plain
terms; a reader interested only in the operator theory can skip it.

A guiding picture is the functional-equation symmetry of the Riemann
zeta function.  The map \(s\mapsto1-\bar s\) is an involution of the
complex plane whose fixed locus is the critical line
\(\operatorname{Re}s=\tfrac12\), and \(\operatorname{Re}s-\tfrac12\)
is the coordinate that changes sign.  A measure on the nontrivial
zeros, weighted by multiplicity, is confined to the fixed locus
exactly when the Riemann hypothesis holds.  The picture explains the
shape of the abstraction; no result about \(\zeta\) is used or
claimed in this paper.  A companion paper \citep{TsiokosRHviaSDTC2026}
studies this instance under explicitly stated hypotheses.

\paragraph{Results.}
Throughout, \((X,J,\mu,\psi,Y,\Jiso)\) is an involutive object ledger
(\Cref{def:duality_confinement:involutive-object-ledger}) with
\(\int_X\|\psim\|^2\,d\mu<\infty\).
\begin{enumerate}
  \item \emph{Trace identity}
    (\Cref{lem:duality_confinement:trace-identity}).  \(\AX\) is a
    positive trace-class operator with
    \(\trace\AX=\int_X\|\psim\|^2\,d\mu\).
  \item \emph{Direct confinement}
    (\Cref{thm:duality_confinement:separation-confinement,thm:duality_confinement:quantitative-confinement}).
    If \(\AX=0\) and \(\psim\) is separating, \(\mu\)-almost every
    point is fixed by \(J\).  If \(X\) is a metric space and
    \(\|\psim(x)\|\geq m(\varepsilon)>0\) whenever
    \(\dist(x,\Fix(J))\geq\varepsilon\), then the set of points at
    distance at least \(\varepsilon\) from \(\Fix(J)\) has outer measure
    at most \(\trace\AX/m(\varepsilon)^2\).
  \item \emph{Douglas factorization}
    (\Cref{thm:duality_confinement:douglas-domination}).  For bounded
    operators \(V,W\) with a common domain,
    \(V^*V\preceq W^*W\) if and only if \(V=TW\) for a contraction
    \(T\).  Applied to the analysis operator of \(\psim\), this
    describes exact domination of \(\AX\) as a factorization problem.
  \item \emph{Confinement theorem}
    (\Cref{thm:duality_confinement:master-theorem}).  If
    \(\AX\preceq B_n\) for positive trace-class \(B_n\) with
    \(\trace B_n\to0\), then \(\AX=0\), \(\psim=0\) almost
    everywhere, and, if \(\psim\) is separating, \(\mu\)-almost every
    point is fixed by \(J\).  Conversely, if \(\psim=0\) almost
    everywhere, the zero sequence is such a family of records.
  \item \emph{Finite windows and defect budgets}
    (\Cref{thm:duality_confinement:exhaustive-squeeze} and
    \Cref{prop:duality_confinement:optimized-trace-budget}).
    The same conclusion holds when the domination is supplied on
    finite windows: real Hilbert spaces \(Y_n\), bounded operators
    \(\iota_n:Y_n\to Y\), positive trace-class operators
    \(A_{X,n}\preceq B_n\) on \(Y_n\), and positive trace-class
    operators \(T_n\) on \(Y\), with
    \(\AX\preceq\iota_nA_{X,n}\iota_n^*+T_n\) and
    \(\trace(\iota_nB_n\iota_n^*)+\trace T_n\to0\).  When \(\AX\) is dominated by
    \((1+t)(K+E_{\mathrm{src}})+(1+t^{-1})E_{\mathrm{br}}\) for every
    \(t>0\), with \(K,E_{\mathrm{src}},E_{\mathrm{br}}\) positive
    trace-class, then
    \(\trace\AX\leq\bigl(\sqrt a+\sqrt b\bigr)^2\) with
    \(a=\trace(K+E_{\mathrm{src}})\), \(b=\trace E_{\mathrm{br}}\), and
    this value is the infimum over \(t>0\) of the scalar budgets.
\end{enumerate}
Every numbered result in this list, and the abstract squeeze of
\Cref{rem:duality_confinement:abstract-cone}, is formalized in Lean~4
with Mathlib for real Hilbert spaces of arbitrary dimension.  The formal
proofs take the trace of a positive operator to be the sum of its
diagonal entries in a Hilbert basis and the Gram operator to be a
Bochner integral in the operator norm.  Bochner integrability in the
trace class, the standard facts identifying the two traces, and the
application of Douglas factorization to the analysis operator are
checked only in the text.  \Cref{app:formalization} lists the
correspondence.

\paragraph{Organization.}
\Cref{sec:framework} explains the framework vocabulary.
\Cref{sec:involutive_ledger,sec:anti_invariant_ledger} define the
ledger, the odd readout, the operator \(\AX\) and separation.
\Cref{sec:direct_confinement} proves the direct confinement results,
\Cref{sec:domination} discusses domination and Douglas factorization,
\Cref{sec:master_theorem} proves the confinement theorem, and
\Cref{sec:exhaustive_squeeze_and_budgets} gives the finite-window and
defect-budget refinements.  \Cref{sec:scope,sec:discussion} discuss
scope and possible instances, and \Cref{sec:conclusion} concludes.
 \section{Framework}
\label{sec:framework}

This section explains the few framework terms used in the paper.
None of them enters a definition, statement or proof in
\Cref{sec:involutive_ledger,sec:anti_invariant_ledger,sec:direct_confinement,sec:domination,sec:master_theorem,sec:exhaustive_squeeze_and_budgets};
a reader interested only in the operator theory may skip to
\Cref{sec:involutive_ledger}.

\paragraph{Closures.}
The author's Six Birds framework \citep{TsiokosFoundationsI2026}
describes a mathematical theory as the collection of fixed points of an
idempotent operator: a completion rule that, applied twice, gives the
same result as applied once.  The fixed points are the objects that the
theory regards as complete, or \emph{closed}.  Later papers add a
discipline for which candidate objects are admissible
\citep{TsiokosFoundationsII2026} and a finite vocabulary for what an
observer can see and certify about them
\citep{TsiokosFoundationsIII2026}.  The title of this paper speaks of
\emph{formed closures}; in the body the phrase is used only informally,
for a concrete structure---here an involutive object ledger---that an
application has specified completely.  No definition or result below
depends on it.

\paragraph{Structural laws and recognition sources.}
The framework distinguishes two kinds of statement.  A \emph{theorem}
is proved from stated hypotheses in the usual way.  A \emph{structural
law} is a named hypothesis asserting that the structures of a
specified class carry certain data.  It is stated separately, with a
name, so that every result that uses it can be seen to use it.  The
data whose existence the law asserts are called its \emph{recognition
source}: the part of an argument that the structure itself must
supply and that is not derived from anything else.

In this paper the structural law is Self-Dual Trace Confinement, whose
name is abbreviated \(\GamSDTC\) or SDTC:
\begin{quote}
  \emph{In a specified class of closures that carry a genuine
  involutive self-duality, observed through a separating odd readout,
  the structure supplies positive trace-class operators \(B_n\) with
  \(\AX\preceq B_n\) and \(\trace B_n\to0\).}
\end{quote}
The recognition source is the sequence \((B_n)\).  The word
``genuine'' is not given a formal meaning, and the class of closures
must be specified in each application; the hypothesis has content
only once both are fixed.  It cannot be asserted for all involutive
ledgers: the complex-conjugation example of \Cref{sec:involutive_ledger}
with Lebesgue measure on a disc centred on the real axis is an
involutive ledger with a
separating odd readout whose mass is not confined, so by
\Cref{thm:duality_confinement:master-theorem} it has no such
sequence.  For any single ledger with a separating odd readout, SDTC is
equivalent to confinement (\Cref{sec:scope}).  Nothing in this paper establishes the hypothesis
in any instance.

\paragraph{What the argument uses.}
The confinement theorem uses the operator \(\AX\) only through its
trace, the Loewner order and the fact that a positive operator of
trace zero is zero.  This is why one statement covers finite weighted
configurations, measure spaces and, in
\Cref{rem:duality_confinement:abstract-cone}, abstract ordered sets
with a faithful trace.
 \section{Involutive object ledger}
\label{sec:involutive_ledger}

The data of the paper are a measured space of objects with an
involution, and an observable compatible with it.

\begin{definition}[Involutive object ledger]
\label{def:duality_confinement:involutive-object-ledger}
An \emph{involutive object ledger} is a tuple
\[
  \IOL=(X,J,\mu,\psi,Y,\Jiso)
\]
consisting of
\begin{enumerate}
  \item a measurable space \(X\), the \emph{carrier};
  \item an involution \(J:X\to X\), that is, \(J(Jx)=x\) for every
    \(x\in X\);
  \item a positive measure \(\mu\) on \(X\);
  \item a real Hilbert space \(Y\), the \emph{response space};
  \item a linear isometric involution \(\Jiso:Y\to Y\), that is,
    \(\Jiso^2=I_Y\) and
    \(\langle\Jiso v,\Jiso w\rangle=\langle v,w\rangle\) for all
    \(v,w\in Y\);
  \item a strongly measurable map \(\psi:X\to Y\), the
    \emph{readout}, that is \emph{equivariant}:
    \[
      \psi(Jx)=\Jiso\psi(x)
      \qquad (x\in X).
    \]
\end{enumerate}
The \emph{fixed locus} of the ledger is
\(\Fix(J)=\{x\in X: Jx=x\}\).
\end{definition}

The carrier-side involution \(J\) acts on the objects whose mass is
being studied, while \(\Jiso\) acts on the space in which those objects
are observed; equivariance says that observing \(Jx\) is the same as
observing \(x\) and then applying \(\Jiso\).  Strong measurability
(being a pointwise limit of measurable simple functions) is the
standard condition under which Bochner integrals of \(Y\)-valued maps
are defined; when \(Y\) is separable it is equivalent to Borel
measurability.  Every result below remains true if \(\psi\) is only
\(\mu\)-almost everywhere strongly measurable (measurability statements
then hold up to \(\mu\)-null sets), and the formal verification uses
that weaker hypothesis.  The measure \(\mu\) is not assumed
\(\sigma\)-finite; under the standing hypothesis
\(\int_X\|\psim\|^2\,d\mu<\infty\) of the later sections, the set
\(\{\psim\neq0\}\) is automatically \(\sigma\)-finite.  No measurability of \(J\) and no invariance of \(\mu\)
under \(J\) is assumed: none of the results below needs them.

\paragraph{The finite weighted case.}
If \(X\) is a finite set with all subsets measurable and every point
has finite mass, the measure is a family of real weights
\(\mu(x)\geq0\), and integrals are finite sums; this is the
\emph{finite weighted case}.  The \emph{visible support} of \(\mu\) is the set of objects of
positive weight.  In general, a statement holds for \(\mu\)-almost
every \(x\) when the set where it fails is contained in a measurable
set of measure zero; in the finite case this means that it holds at
every visible object.  The mass of the ledger is \emph{confined to the
fixed locus} when \(Jx=x\) for \(\mu\)-almost every \(x\), that is,
when \(X\setminus\Fix(J)\) is \(\mu\)-null.  If \(\Fix(J)\) is
measurable this says \(\mu(X\setminus\Fix(J))=0\).

\begin{example}[Complex conjugation]
Let \(X\subset\mathbb{C}\) be a Borel set stable under complex
conjugation, with \(\mu\) any positive Borel measure on \(X\).  Put
\[
  J(z)=\overline{z},\qquad
  Y=\mathbb{R}^2,\qquad
  \Jiso(u,v)=(u,-v),\qquad
  \psi(z)=(\operatorname{Re}z,\operatorname{Im}z).
\]
Then \(\psi(Jz)=(\operatorname{Re}z,-\operatorname{Im}z)=\Jiso\psi(z)\),
so the readout is equivariant.  The fixed locus is \(X\cap\mathbb{R}\),
and the coordinate that changes sign is the imaginary part.  Mass is
confined to the fixed locus exactly when \(\mu\) is carried by the real
axis.
\end{example}

\begin{example}[Zeros of the Riemann zeta function]
The same pattern underlies the motivating example.  The map
\(J_L(s)=1-\overline{s}\) is an involution of \(\mathbb{C}\) with fixed
locus the critical line \(\operatorname{Re}s=\tfrac12\).  Taking for
\(X\) the nontrivial zeros of the Riemann zeta function, \(\mu\) the
counting measure weighted by multiplicity (the zeros are symmetric
under \(J_L\) by the functional equation and the reflection principle
\citep{TitchmarshHeathBrown1986}), \(Y=\mathbb{R}\),
\(\Jiso=-I\) and \(\psi(s)=\operatorname{Re}s-\tfrac12\), one obtains an
involutive object ledger whose mass is confined to the fixed locus
exactly when the Riemann hypothesis holds.  Here \(\psi\) is already
odd, so \(\psim=\psi\).  This example only illustrates the definition.
The results below also require
\(\int_X\|\psim\|^2\,d\mu=\sum_\rho(\operatorname{Re}\rho-\tfrac12)^2<\infty\),
summed over the zeros with multiplicity; the sum vanishes if the
Riemann hypothesis holds, but its finiteness is not known otherwise.
Nothing in this paper bears on the location of the zeros.
\end{example}
 \section{Anti-invariant ledger and separation}
\label{sec:anti_invariant_ledger}

Fix an involutive object ledger \(\IOL=(X,J,\mu,\psi,Y,\Jiso)\).  This
section extracts the odd part of the readout, builds the positive
operator \(\AX\) from it, and computes its trace.

\paragraph{The odd readout.}
Because \(\Jiso\) is a linear isometric involution it is self-adjoint:
\(\Jiso^*=\Jiso^{-1}=\Jiso\).  Hence
\[
  \Pminus=\tfrac12\,(I_Y-\Jiso)
\]
satisfies \(\Pminus^2=\Pminus=\Pminus^*\); it is the orthogonal
projection onto the \((-1)\)-eigenspace \(\{v\in Y:\Jiso v=-v\}\).
The \emph{anti-invariant} (or odd) readout is
\[
  \psim(x)=\Pminus\psi(x)\qquad(x\in X).
\]
It is strongly measurable, and by equivariance
\[
  \psim(Jx)=\Pminus\Jiso\psi(x)=-\Pminus\psi(x)=-\psim(x),
\]
because \(\Pminus\Jiso=-\Pminus\).  In particular, if \(Jx=x\) then
\(\psim(x)=-\psim(x)\), so \(\psim(x)=0\): the odd readout vanishes on
the fixed locus.  The converse is the separation property defined
below.  Only this vanishing is used later: none of the results needs
\(J\) to be an involution, and all of them hold for any map
\(J:X\to X\) with \(\psi\circ J=\Jiso\circ\psi\).  The involutive
setting is the one in which the fixed locus has the meaning of a
self-dual locus, and we keep it for that reason.

\paragraph{Positive and trace-class operators.}
On a real Hilbert space an operator \(T\) is \emph{positive}, written
\(T\succeq0\), if it is self-adjoint and \(\langle Th,h\rangle\geq0\)
for all \(h\); the Loewner order is \(A\preceq B\) if and only if
\(B-A\succeq0\).  For \(v\in Y\) write \(vv^*\) for the rank-one
operator \(h\mapsto\langle v,h\rangle v\); it is positive.  Let
\(\mathcal S_1(Y)\) denote the trace class with its norm
\(\|\cdot\|_1\).  We use four standard facts
\citep{Simon2005,ReedSimon1980}:
\begin{enumerate}[label=(T\arabic*)]
  \item \(\mathcal S_1(Y)\) is a Banach space, and the trace is a
    continuous linear functional with \(|\trace T|\leq\|T\|_1\); for
    \(v\in Y\), \(\|vv^*\|_1=\trace(vv^*)=\|v\|^2\).
  \item For \(h,k\in Y\), the map \(T\mapsto\langle Th,k\rangle\) is a
    continuous linear functional on \(\mathcal S_1(Y)\).
  \item If \(T\succeq0\) is trace-class, then
    \(\trace T=\sum_i\langle Te_i,e_i\rangle\geq0\) for every
    orthonormal basis \((e_i)\), and \(\trace T=0\) implies \(T=0\).
  \item If \(0\preceq A\preceq B\) and \(B\) is trace-class, then \(A\)
    is trace-class and \(\trace A\leq\trace B\).
\end{enumerate}
These are usually stated for complex Hilbert spaces; for a real
Hilbert space they follow by passing to the complexification
\(Y\otimes_{\mathbb R}\mathbb C\), which preserves positivity of
self-adjoint operators, the trace norm and the trace.  When \(Y\) is
finite-dimensional every operator is trace-class and (T1)--(T4) are
facts of linear algebra.  The formal verification
(\Cref{app:formalization}) does not use (T1)--(T4): it proves the
results with the trace of a positive operator defined as
\(\sum_i\langle Te_i,e_i\rangle\in[0,\infty]\).  By (T3) this agrees
with the trace on positive trace-class operators, and a positive
operator for which it is finite is trace-class
\citep{ReedSimon1980}.

\begin{definition}[Anti-invariant object ledger]
\label{def:duality_confinement:anti-invariant-ledger}
Assume \(\int_X\|\psim(x)\|^2\,d\mu(x)<\infty\).  The
\emph{anti-invariant object ledger} of \(\IOL\) is the operator
\[
  \AX=\int_X \psim(x)\,\psim(x)^*\,d\mu(x),
\]
the integral being a Bochner integral in \(\mathcal S_1(Y)\).
\end{definition}

The integral exists by \Cref{lem:duality_confinement:trace-identity}
below.  For a finite weighted ledger it is the finite sum
\[
  \AX=\sum_{x\in X}\mu(x)\,\psim(x)\psim(x)^*,
\]
a nonnegative combination of rank-one positive operators, one for each
visible object, each recording the part of the readout of that object
that changes sign under the duality.

\begin{definition}[Separating anti-invariant readout]
\label{def:duality_confinement:separating-readout}
The odd readout \(\psim\) is \emph{qualitatively separating} if, for
\(\mu\)-almost every \(x\in X\),
\[
  \psim(x)=0
  \quad\Longleftrightarrow\quad
  x\in\Fix(J).
\]
In the finite weighted case this means that the equivalence holds at
every object of positive weight.

Suppose in addition that \(X\) is a metric space.  The odd readout is
\emph{quantitatively separating} with modulus
\(m:(0,\infty)\to(0,\infty)\) if, for every \(\varepsilon>0\) and every
\(x\in X\),
\[
  \dist(x,\Fix(J))\geq\varepsilon
  \quad\Longrightarrow\quad
  \|\psim(x)\|\geq m(\varepsilon).
\]
\end{definition}

Since \(\psim\) always vanishes on \(\Fix(J)\), qualitative separation
only adds the implication \(\psim(x)=0\Rightarrow Jx=x\).  Quantitative
separation says that points at distance at least \(\varepsilon\) from the
fixed locus have odd readout at least \(m(\varepsilon)\) in norm
(with the convention \(\dist(x,\varnothing)=+\infty\)).  It implies
that \(\psim(x)=0\) only if \(x\) lies in the closure of \(\Fix(J)\), so
it implies qualitative separation when \(\Fix(J)\) is closed, for
instance when \(J\) is continuous.  The metric enters only through the function \(x\mapsto\dist(x,\Fix(J))\);
in \Cref{thm:duality_confinement:quantitative-confinement} it may be
replaced by any function \(\delta:X\to[0,\infty]\) for which the
displayed implication holds.

\begin{lemma}[Trace identity]
\label{lem:duality_confinement:trace-identity}
Let \(\IOL\) be an involutive object ledger with
\(\int_X\|\psim\|^2\,d\mu<\infty\).  Then the map
\(x\mapsto\psim(x)\psim(x)^*\) is Bochner integrable in
\(\mathcal S_1(Y)\), the operator \(\AX\) is positive and trace-class,
\[
  \langle\AX h,k\rangle=\int_X\langle\psim(x),h\rangle\,
  \langle\psim(x),k\rangle\,d\mu(x)\qquad(h,k\in Y),
\]
and
\[
  \trace\AX=\int_X\|\psim(x)\|^2\,d\mu(x).
\]
\end{lemma}

\begin{proof}
For \(u,v\in Y\) we have \(uu^*-vv^*=u(u-v)^*+(u-v)v^*\), and a
rank-one operator \(ab^*\) has trace norm \(\|a\|\,\|b\|\).  Hence
\(\|uu^*-vv^*\|_1\leq(\|u\|+\|v\|)\,\|u-v\|\), so \(v\mapsto vv^*\) is
continuous from \(Y\) to \(\mathcal S_1(Y)\).  Composed with the
strongly measurable map \(\psim\), it gives a strongly measurable map
\(x\mapsto\psim(x)\psim(x)^*\) whose norm
\(\|\psim(x)\psim(x)^*\|_1=\|\psim(x)\|^2\) is integrable by (T1).
The map is therefore Bochner integrable and \(\AX\in\mathcal S_1(Y)\).

Continuous linear functionals commute with Bochner integrals.  Applying
this to the functionals of (T2) gives the displayed formula for
\(\langle\AX h,k\rangle\).  Its right-hand side is symmetric in \(h\)
and \(k\), and for \(h=k\) it equals
\(\int_X\langle\psim,h\rangle^2\,d\mu\geq0\); so \(\AX\) is positive.
Applying the same principle to the trace, which is a continuous linear
functional by (T1), and using \(\trace(vv^*)=\|v\|^2\), gives
\(\trace\AX=\int_X\|\psim\|^2\,d\mu\).
\end{proof}

By (T3) and the trace identity, the following are equivalent:
\(\AX=0\); \(\trace\AX=0\); \(\psim(x)=0\) for \(\mu\)-almost every
\(x\).  The last equivalence holds because the integral of the
nonnegative function \(\|\psim\|^2\) vanishes exactly when the function
vanishes almost everywhere.

\begin{remark}[Analysis operator]
\label{rem:duality_confinement:analysis-operator}
Under the hypothesis of \Cref{lem:duality_confinement:trace-identity},
define \(V:Y\to L^2(X,\mu)\) by \((Vh)(x)=\langle\psim(x),h\rangle\).
By the Cauchy--Schwarz inequality
\(\|Vh\|^2\leq\|h\|^2\int_X\|\psim\|^2\,d\mu\), so \(V\) is bounded, and
by the lemma \(\langle\AX h,h\rangle=\|Vh\|^2=\langle V^*Vh,h\rangle\).
Two self-adjoint operators with the same quadratic form are equal, so
\[
  \AX=V^*V.
\]
In particular \(V\) is a Hilbert--Schmidt operator with
\(\|V\|_{\mathrm{HS}}^2=\trace V^*V=\trace\AX=\int_X\|\psim\|^2\,d\mu\).
Thus \(\AX\) is the Gram operator of the odd readout, and domination
questions about \(\AX\) are questions about \(V\); see
\Cref{sec:domination}.
\end{remark}
 \section{Direct confinement}
\label{sec:direct_confinement}

The trace identity turns information about \(\AX\) into information
about where the mass of \(\mu\) lies.  The qualitative statement uses
\(\AX=0\); the quantitative one uses only the size of \(\trace\AX\).

\begin{theorem}[Separation--confinement]
\label{thm:duality_confinement:separation-confinement}
Let \(\IOL\) be an involutive object ledger with
\(\int_X\|\psim\|^2\,d\mu<\infty\), and suppose \(\AX=0\).  Then
\(\psim(x)=0\) for \(\mu\)-almost every \(x\).  If moreover \(\psim\)
is qualitatively separating, then \(Jx=x\) for \(\mu\)-almost every
\(x\); that is, \(X\setminus\Fix(J)\) is \(\mu\)-null.  In the finite
weighted case, every object of positive weight is fixed by \(J\).
\end{theorem}

\begin{proof}
By \Cref{lem:duality_confinement:trace-identity},
\(\int_X\|\psim\|^2\,d\mu=\trace\AX=0\).  The integrand is nonnegative
and measurable, so \(\|\psim(x)\|=0\) for \(\mu\)-almost every \(x\).
If \(\psim\) is qualitatively separating, the set of \(x\) at which the
separation equivalence fails is also \(\mu\)-null; outside the union of
the two null sets, \(\psim(x)=0\) and hence \(Jx=x\).  In the finite
case the null sets contain no object of positive weight.
\end{proof}

\begin{theorem}[Quantitative confinement]
\label{thm:duality_confinement:quantitative-confinement}
Let \(\IOL\) be an involutive object ledger on a metric space \(X\),
with \(\int_X\|\psim\|^2\,d\mu<\infty\) and with \(\psim\)
quantitatively separating with modulus \(m\).  Then, for every
\(\varepsilon>0\),
\[
  \mu^*\bigl(\{x\in X:\dist(x,\Fix(J))\geq\varepsilon\}\bigr)
  \leq
  \frac{\trace\AX}{m(\varepsilon)^2},
\]
where \(\mu^*\) is the outer measure of \(\mu\); it equals \(\mu\) on
measurable sets.  The same bound holds with \(\dist(\cdot,\Fix(J))\)
replaced by any function \(\delta:X\to[0,\infty]\) such that
\(\delta(x)\geq\varepsilon\) implies \(\|\psim(x)\|\geq m(\varepsilon)\).
\end{theorem}

\begin{proof}
Fix \(\varepsilon>0\) and put
\(E_\varepsilon=\{x:\dist(x,\Fix(J))\geq\varepsilon\}\) and
\(F_\varepsilon=\{x:\|\psim(x)\|\geq m(\varepsilon)\}\).  Quantitative
separation gives \(E_\varepsilon\subseteq F_\varepsilon\), and
\(F_\varepsilon\) is measurable because \(\psim\) is strongly
measurable (up to a \(\mu\)-null set if \(\psi\) is only almost
everywhere strongly measurable, which does not affect the bound).  Markov's inequality \citep{Folland1999} and
the trace identity give
\[
  m(\varepsilon)^2\,\mu(F_\varepsilon)
  \leq\int_{F_\varepsilon}\|\psim\|^2\,d\mu
  \leq\int_X\|\psim\|^2\,d\mu=\trace\AX .
\]
Since \(m(\varepsilon)>0\), dividing and using
\(\mu^*(E_\varepsilon)\leq\mu(F_\varepsilon)\) gives the bound.  The
proof used the metric only through the inclusion
\(E_\varepsilon\subseteq F_\varepsilon\), which holds for any
\(\delta\) as in the statement.
\end{proof}

\begin{example}[Complex conjugation, continued]
In the example of \Cref{sec:involutive_ledger},
\(\Pminus(u,v)=(0,v)\), so \(\psim(z)=(0,\operatorname{Im}z)\) and
\[
  \AX=\Bigl(\int_X(\operatorname{Im}z)^2\,d\mu(z)\Bigr)
  \begin{pmatrix}0&0\\0&1\end{pmatrix},
\]
provided the integral is finite.  If \(X\) is a disc centred on the
real axis, then \(\operatorname{Re}z\in X\) for every \(z\in X\), so
\(\dist(z,X\cap\mathbb R)=|\operatorname{Im}z|\); the readout is then
quantitatively separating with \(m(\varepsilon)=\varepsilon\), and
\Cref{thm:duality_confinement:quantitative-confinement} becomes the
familiar Markov--Chebyshev bound
\(\mu\{|\operatorname{Im}z|\geq\varepsilon\}\leq
\varepsilon^{-2}\int_X(\operatorname{Im}z)^2\,d\mu\).
\end{example}

The confinement theorem of \Cref{sec:master_theorem} supplies the
missing input \(\AX=0\): it derives it from domination by operators of
vanishing trace, after which
\Cref{thm:duality_confinement:separation-confinement} applies.  The
quantitative theorem is useful when only an upper bound on
\(\trace\AX\) is available, as in
\Cref{prop:duality_confinement:optimized-trace-budget}.
 \section{Domination and Douglas factorization}
\label{sec:domination}

\Cref{sec:direct_confinement} shows what follows once \(\AX\) vanishes.
To force the vanishing, the confinement theorem uses upper bounds on
\(\AX\) in the Loewner order.  In applications such bounds typically
come in two pieces: a main term, which an argument specific to the
application makes small, and an error term, kept explicit.

\begin{definition}[Completed domination bridge]
\label{def:duality_confinement:completed-domination-bridge}
Let \(\IOL\) be an involutive object ledger with
\(\int_X\|\psim\|^2\,d\mu<\infty\).  A \emph{completed domination
bridge} for \(\AX\) is a pair of positive trace-class operators
\(\Kminus\) and \(E\) on \(Y\) such that
\[
  \AX\preceq\Kminus+E.
\]
The operator \(\Kminus\) is the \emph{main term} and \(E\) the
\emph{defect}.  The bridge is \emph{exact} when \(E=0\).  A sequence of
bridges \(\AX\preceq\Kminus_n+E_n\) is a \emph{domination-record
sequence} if the bounds \(B_n=\Kminus_n+E_n\) satisfy
\(\trace B_n\to0\) as \(n\to\infty\).
\end{definition}

A bridge is bookkeeping rather than a theorem: it records which part of
an upper bound is the controlled main term and which part is error.
As a definition it imposes nothing beyond \(\AX\preceq\Kminus+E\);
indeed \(\Kminus=\AX\), \(E=0\) is always an exact bridge.  Its content
in applications is that \(\Kminus\) is computed from data that do not
involve \(\psim\).
The notation follows the framework's companion papers, where the main
term is computed on the carrier side from other data and the bridge
inequality is established by a separate argument.  For example, in the
finite-dimensional setting of \citep{TsiokosNeedleKiller2026},
\(\Kminus=L^-C^\dagger(L^-)^*\), where \(C\) is a positive audit
operator on a carrier space, \(C^\dagger\) its Moore--Penrose
pseudoinverse, and \(L^-\) an anti-invariant probe with
\(\ker C\subseteq\ker L^-\); the minus sign records that \(\Kminus\) is
the currency of the anti-invariant probe.  A
domination-record sequence is exactly the input of the confinement
theorem (\Cref{thm:duality_confinement:master-theorem}).  Exact
bridges have a useful reformulation, which rests on the following
classical theorem of Douglas.

\begin{theorem}[Douglas factorization {\citep{Douglas1966}}]
\label{thm:duality_confinement:douglas-domination}
Let \(H\), \(K_1\) and \(K_2\) be real Hilbert spaces and let
\(V:H\to K_1\) and \(W:H\to K_2\) be bounded linear operators.  Then
\[
  V^*V\preceq W^*W
  \quad\Longleftrightarrow\quad
  V=TW\ \text{for some bounded } T:K_2\to K_1 \text{ with } \|T\|\leq1.
\]
\end{theorem}

\begin{proof}
If \(V=TW\) with \(\|T\|\leq1\), then for every \(h\in H\),
\[
  \langle V^*Vh,h\rangle=\|TWh\|^2\leq\|Wh\|^2=\langle W^*Wh,h\rangle,
\]
so \(V^*V\preceq W^*W\).  Conversely, suppose \(V^*V\preceq W^*W\), so
that \(\|Vh\|\leq\|Wh\|\) for every \(h\in H\).  Define \(T_0\) on the
range of \(W\) by \(T_0(Wh)=Vh\).  It is well defined and linear: if
\(Wh=Wh'\) then \(\|Vh-Vh'\|\leq\|W(h-h')\|=0\).  It is contractive by
the same inequality.  Hence \(T_0\) extends by continuity to a
contraction on the closure \(\overline{\operatorname{Ran}W}\).  Setting
\(T=T_0\) on \(\overline{\operatorname{Ran}W}\) and \(T=0\) on its
orthogonal complement gives a bounded operator \(T:K_2\to K_1\) with
\(\|T\|\leq1\) and \(TW=V\).
\end{proof}

The theorem is due to Douglas \citep{Douglas1966}, who proves, for
bounded operators \(A,B\) on a single complex Hilbert space, that
\(AA^*\preceq\lambda^2BB^*\) for some \(\lambda\geq0\) if and only if
\(A=BC\) for a bounded operator \(C\), and that \(C\) can then be
chosen with \(\|C\|\leq\lambda\).  With \(\lambda=1\), \(A=V^*\) and
\(B=W^*\) this is the displayed statement after taking adjoints.  The
short proof above, which is standard, covers real scalars and
operators between different spaces.

\paragraph{Exact bridges as factorizations.}
By \Cref{rem:duality_confinement:analysis-operator},
\(\AX=V^*V\) for the analysis operator
\(V:Y\to L^2(X,\mu)\), \((Vh)(x)=\langle\psim(x),h\rangle\).  Every
positive operator \(\Kminus\) on \(Y\) can be written as
\(\Kminus=W^*W\), for instance with \(W=(\Kminus)^{1/2}\), which is
Hilbert--Schmidt when \(\Kminus\) is trace-class.
\Cref{thm:duality_confinement:douglas-domination} therefore says that
\(\AX\preceq\Kminus\) exactly when the analysis operator of the odd
readout factors as \(V=TW\) through a contraction \(T\).  In words: an
exact bridge exists precisely when every odd measurement
\(x\mapsto\langle\psim(x),h\rangle\) is obtained from the main-term
measurement \(Wh\) by one fixed contraction.
 \section{The duality-confinement membrane theorem}
\label{sec:master_theorem}

This section proves the confinement theorem.  Its input is a sequence
of domination records whose traces vanish; its output is the
almost-everywhere vanishing of the odd readout and, when the readout is
separating, confinement of the mass of \(\mu\) to the fixed locus.  In the
framework's companion papers a positive budget of this kind, bounding
a carrier-side quantity at every stage, is called a \emph{membrane}
\citep{TsiokosNeedleKiller2026}; the theorem's name records that it is
the step turning such budgets into confinement.  In the vocabulary of
\Cref{sec:framework}, it is the mathematical consequence of the
structural law \(\GamSDTC\): it does not produce the domination
records, it uses them.

\begin{theorem}[Duality-confinement membrane theorem]
\label{thm:duality_confinement:master-theorem}
Let \(\IOL=(X,J,\mu,\psi,Y,\Jiso)\) be an involutive object ledger with
\(\int_X\|\psim\|^2\,d\mu<\infty\).  Suppose there are positive
trace-class operators \(B_n\) on \(Y\) such that
\[
  \AX\preceq B_n\quad(n=1,2,\ldots)
  \qquad\text{and}\qquad
  \trace B_n\to0\quad(n\to\infty).
\]
Then \(\AX=0\) and \(\psim(x)=0\) for \(\mu\)-almost every \(x\).  If,
in addition, \(\psim\) is qualitatively separating, then \(Jx=x\) for
\(\mu\)-almost every \(x\), so the mass of \(\mu\) is confined to
\(\Fix(J)\); in the finite weighted case, every object of positive
weight is fixed by \(J\).

Conversely, if \(\psim(x)=0\) for \(\mu\)-almost every \(x\)---in
particular, if \(Jx=x\) for \(\mu\)-almost every \(x\)---then
\(\AX=0\), and the constant sequence \(B_n=0\) satisfies the
hypotheses.
\end{theorem}

\begin{proof}
By \Cref{lem:duality_confinement:trace-identity}, \(\AX\) is positive
and trace-class.  For each \(n\), (T4) applied to
\(0\preceq\AX\preceq B_n\) gives
\[
  0\leq\trace\AX\leq\trace B_n .
\]
The middle quantity does not depend on \(n\) and the right-hand side
tends to zero, so \(\trace\AX=0\).  By (T3), \(\AX=0\).  The remaining
direct assertions are
\Cref{thm:duality_confinement:separation-confinement}.  For the
converse, \(\psim\) vanishes on \(\Fix(J)\), so almost-everywhere
fixedness implies \(\psim=0\) almost everywhere; then
\(\trace\AX=\int_X\|\psim\|^2\,d\mu=0\), so \(\AX=0\) by (T3), and
\(\AX\preceq0\).
\end{proof}

The converse shows that, for a given ledger, trace-vanishing domination
records exist if and only if \(\psim\) vanishes almost everywhere, and
hence, for a separating readout, if and only if the mass is confined to
the fixed locus.  The hypothesis is therefore equivalent to the
conclusion, and the theorem does not make confinement easier to prove.
Its role is to fix the form in which confinement is to be established
when \(\AX\) is known only through a bridge
\(\AX\preceq\Kminus+E\) to operators computed from other data
(\Cref{sec:domination}): then trace bounds on \(\Kminus\) and \(E\),
rather than on \(\AX\) itself, are what must be proved.

Two features of the proof are worth isolating.  First, it uses \(\AX\)
only through four properties: \(\AX\) is positive; the trace is
monotone for the Loewner order; the trace of a positive operator is
nonnegative; and a positive operator of trace zero is zero.  Second,
the geometric conclusion enters only at the end, through the trace
identity and separation.  The first observation gives an abstract form
of the theorem.

\begin{remark}[Abstract form]
\label{rem:duality_confinement:abstract-cone}
Let \(C\) be a set with a distinguished element \(0\), a relation
\(\preceq\), a subset of \emph{positive} elements, and a map
\(\tau:C\to\mathbb R\) such that \(\tau(c)\geq0\) for every positive
\(c\), and \(a\preceq b\) implies \(\tau(a)\leq\tau(b)\).
Let \(a\in C\) be positive with \(\tau(a)=0\) only if \(a=0\).  If
\(a\preceq b_n\) for all \(n\) and \(\tau(b_n)\to0\), then \(a=0\): indeed \(0\leq\tau(a)\leq\tau(b_n)\to0\).  \Cref{thm:duality_confinement:master-theorem} is
the case of positive trace-class operators with the Loewner order and
the operator trace.
\end{remark}

\paragraph{Where the theorem is used.}
A companion paper \citep{TsiokosRHviaSDTC2026} considers the
involution \(s\mapsto1-\bar s\) on the zeros of the Riemann zeta
function.  It states, as an explicit hypothesis, a recognition source
supplying domination records for an odd readout of the zeros, and
applies the fixed-locus conclusion of the theorem to that ledger.  The
present paper uses nothing from that work.
 \section{Exhaustive squeeze and defected budgets}
\label{sec:exhaustive_squeeze_and_budgets}

In applications the domination records are rarely available on \(Y\)
itself.  More often one controls a finite piece of the problem---a
window, a truncation, a finite-dimensional subspace---and must account
separately for what the piece misses.  This section records the two
refinements of the confinement theorem needed in that situation.
Throughout, \(\IOL\) is an involutive object ledger with
\(\int_X\|\psim\|^2\,d\mu<\infty\).

\begin{definition}[Exhaustive moving ledger]
\label{def:duality_confinement:exhaustive-moving-ledger}
An \emph{exhaustive moving ledger} for \(\AX\) consists of, for each
\(n\geq1\), a real Hilbert space \(Y_n\), a bounded operator
\(\iota_n:Y_n\to Y\), a positive trace-class operator \(A_{X,n}\) on
\(Y_n\) and a positive trace-class operator \(T_n\) on \(Y\), such that
\[
  \AX\preceq\iota_nA_{X,n}\iota_n^*+T_n\qquad(n\geq1).
\]
It is \emph{vanishing-exhaustive} if there are positive trace-class
operators \(B_n\) on \(Y_n\) with \(A_{X,n}\preceq B_n\) and
\[
  \trace(\iota_nB_n\iota_n^*)+\trace T_n\to0\qquad(n\to\infty).
\]
\end{definition}

The operator \(A_{X,n}\) is the part of the ledger seen by the
\(n\)-th window, \(\iota_n\) transports it into \(Y\), and \(T_n\) is a
positive tail accounting for what the window misses.  Despite the
name, no exhaustion, monotonicity or convergence of the windows is
required; the name reflects the typical use, illustrated next.

\begin{theorem}[Exhaustive ledger squeeze]
\label{thm:duality_confinement:exhaustive-squeeze}
If \(\AX\) admits a vanishing-exhaustive moving ledger, then
\(\AX=0\) and \(\psim(x)=0\) for \(\mu\)-almost every \(x\).  If, in
addition, \(\psim\) is qualitatively separating, then \(Jx=x\) for
\(\mu\)-almost every \(x\).
\end{theorem}

\begin{proof}
Conjugation preserves the Loewner order: if \(A\preceq B\) on \(Y_n\),
then for \(h\in Y\),
\(\langle\iota_n(B-A)\iota_n^*h,h\rangle
=\langle(B-A)\iota_n^*h,\iota_n^*h\rangle\geq0\), so
\(\iota_nA\iota_n^*\preceq\iota_nB\iota_n^*\).  Hence
\[
  \AX\preceq\iota_nA_{X,n}\iota_n^*+T_n
  \preceq\iota_nB_n\iota_n^*+T_n=:B_n' .
\]
Each \(B_n'\) is positive and trace-class, because the trace class is
an ideal, and \(\trace B_n'\to0\) by hypothesis.
\Cref{thm:duality_confinement:master-theorem} applied to the sequence
\((B_n')\) gives the conclusions.
\end{proof}

\Cref{thm:duality_confinement:exhaustive-squeeze} is
\Cref{thm:duality_confinement:master-theorem} applied to the records
\(B_n'\); its content is bookkeeping, recording the data an
application typically computes.

\begin{example}[Finite-dimensional windows]
Let \(Y\) be separable, and let \(P_n\) be orthogonal projections onto
finite-dimensional subspaces \(Y_n\subseteq Y\) with \(P_n\to I\)
strongly, let \(\iota_n\) be the
inclusion of \(Y_n\), so that \(\iota_n\iota_n^*=P_n\), and put
\(Q_n=I-P_n\).  With \(\AX=V^*V\) as in
\Cref{rem:duality_confinement:analysis-operator}, the case \(t=1\) of
the operator inequality displayed below, applied to
\(V=VP_n+VQ_n\), gives
\[
  \AX\preceq2P_n\AX P_n+2Q_n\AX Q_n .
\]
So the compressions \(A_{X,n}=2\,\iota_n^*\AX\iota_n\) and the tails
\(T_n=2Q_n\AX Q_n\) form an exhaustive moving ledger.  The tails
satisfy \(\trace T_n=2\|VQ_n\|_{\mathrm{HS}}^2\to0\) by dominated
convergence, because \(V\) is Hilbert--Schmidt.  Whether the window
bounds \(B_n\) can be chosen with \(\trace(\iota_nB_n\iota_n^*)\to0\) is
the substantive question in any application.
\end{example}

The second refinement concerns a bound with two error terms of
different origin.  Such bounds come from the operator form of the
elementary inequality \((\alpha+\beta)^2\leq(1+t)\alpha^2+(1+t^{-1})\beta^2\):
for bounded operators \(V_1,V_2\) and \(t>0\),
\[
  (1+t)V_1^*V_1+(1+t^{-1})V_2^*V_2-(V_1+V_2)^*(V_1+V_2)
  =\bigl(\sqrt tV_1-V_2/\sqrt t\bigr)^*\bigl(\sqrt tV_1-V_2/\sqrt t\bigr)
  \succeq0 .
\]
So if the analysis operator of \Cref{rem:duality_confinement:analysis-operator}
splits as \(V=V_1+V_2\), then \(\AX\) is dominated, for every \(t>0\),
by \((1+t)V_1^*V_1+(1+t^{-1})V_2^*V_2\).

\begin{definition}[Defected obstruction budget]
\label{def:duality_confinement:defected-budget}
A \emph{defected obstruction budget} for \(\AX\) consists of positive
trace-class operators \(\Kminus\), \(E_{\mathrm{src}}\) and
\(E_{\mathrm{br}}\) on \(Y\) such that
\[
  \AX
  \preceq
  (1+t)\bigl(\Kminus+E_{\mathrm{src}}\bigr)+(1+t^{-1})E_{\mathrm{br}}
  \qquad\text{for every }t>0 .
\]
Here \(E_{\mathrm{src}}\) is a source defect and \(E_{\mathrm{br}}\) a
bridge defect.  A \emph{sequence} of defected obstruction budgets is a
sequence \((\Kminus_n,E_{\mathrm{src},n},E_{\mathrm{br},n})\), each a
defected obstruction budget for the same operator \(\AX\).
\end{definition}

The operators \(\Kminus\) and \(E_{\mathrm{src}}\) enter only through
their sum; they are kept apart because in applications they have
different origins, a computed main term and a defect in its source.

The choice \(t=1\) gives the unweighted bound
\(2(\Kminus+E_{\mathrm{src}})+2E_{\mathrm{br}}\); the next proposition
optimizes over \(t\).

\begin{proposition}[Optimized scalar trace budget]
\label{prop:duality_confinement:optimized-trace-budget}
Let \((\Kminus,E_{\mathrm{src}},E_{\mathrm{br}})\) be a defected
obstruction budget for \(\AX\), and put
\(a=\trace(\Kminus+E_{\mathrm{src}})\) and \(b=\trace E_{\mathrm{br}}\).
Then \(a,b\geq0\),
\[
  \inf_{t>0}\bigl[(1+t)a+(1+t^{-1})b\bigr]
  =\bigl(\sqrt a+\sqrt b\bigr)^2,
\]
the infimum being attained at \(t=\sqrt{b/a}\) when \(a,b>0\), and
\[
  \trace\AX\leq\bigl(\sqrt a+\sqrt b\bigr)^2 .
\]
Consequently, if \((\Kminus_n,E_{\mathrm{src},n},E_{\mathrm{br},n})\)
is a sequence of defected obstruction budgets for \(\AX\) with
\(a_n=\trace(\Kminus_n+E_{\mathrm{src},n})\to0\) and
\(b_n=\trace E_{\mathrm{br},n}\to0\), then all conclusions of
\Cref{thm:duality_confinement:master-theorem} hold.
\end{proposition}

\begin{proof}
The traces \(a\) and \(b\) are nonnegative by (T3).  For \(t>0\),
\[
  (1+t)a+(1+t^{-1})b-\bigl(\sqrt a+\sqrt b\bigr)^2
  =ta+\frac bt-2\sqrt{ab}
  =\Bigl(\sqrt{ta}-\sqrt{b/t}\Bigr)^2\geq0,
\]
with equality when \(ta=b/t\), that is, at \(t=\sqrt{b/a}\) if
\(a,b>0\).  If \(a=0\) the budget equals \(b+b/t\), which decreases to
\(b=(\sqrt a+\sqrt b)^2\) as \(t\to\infty\); if \(b=0\) it equals
\(a+ta\), which decreases to \(a\) as \(t\to0\).  This proves the
formula for the infimum.  For each \(t>0\), (T4) and linearity of the
trace give \(\trace\AX\leq(1+t)a+(1+t^{-1})b\); taking the infimum
gives \(\trace\AX\leq(\sqrt a+\sqrt b)^2\).  For a sequence of budgets,
\(0\leq\trace\AX\leq(\sqrt{a_n}+\sqrt{b_n})^2\to0\), so
\(\trace\AX=0\), hence \(\AX=0\) by (T3), and the remaining
conclusions follow from
\Cref{thm:duality_confinement:separation-confinement}.
\end{proof}

When \(\AX=V^*V\) and \(V=V_1+V_2\) with
\(\Kminus+E_{\mathrm{src}}=V_1^*V_1\) and \(E_{\mathrm{br}}=V_2^*V_2\),
as in the operator inequality above, the bound
\(\trace\AX\leq(\sqrt a+\sqrt b)^2\) is the triangle inequality
\(\|V\|_{\mathrm{HS}}\leq\|V_1\|_{\mathrm{HS}}+\|V_2\|_{\mathrm{HS}}\),
and the optimization over \(t\) is a standard proof of it.  The
proposition applies more generally, to any family of bounds of the
displayed form.

For a single budget, combining the proposition with
\Cref{thm:duality_confinement:quantitative-confinement} gives the
explicit bound
\(\mu^*\{x:\dist(x,\Fix(J))\geq\varepsilon\}\leq
(\sqrt a+\sqrt b)^2/m(\varepsilon)^2\) whenever \(\psim\) is
quantitatively separating with modulus \(m\).
 \section{Scope and non-claims}
\label{sec:scope}

This section states what the paper does not claim.

\paragraph{SDTC is a hypothesis.}
Self-Dual Trace Confinement is a named hypothesis about closures with
genuine involutive self-duality (\Cref{sec:framework}).  It is not
derived from anything in this paper, and it is not established in any
particular instance.  By the converse part of
\Cref{thm:duality_confinement:master-theorem}, in a given ledger the
existence of trace-vanishing domination records is equivalent to the
almost-everywhere vanishing of the odd readout, and hence, when the
odd readout is separating, to confinement.  So SDTC, in any instance
with a separating readout, asserts as much as the confinement
statement it is used to prove; its role is to name the form in which the confinement is to be
established.

\paragraph{Infinite-dimensional response spaces.}
All results are proved in the text for an arbitrary real Hilbert space
\(Y\), using the standard trace-class facts (T1)--(T4) of
\Cref{sec:anti_invariant_ledger}.  The formal verification also covers
arbitrary dimension, with the trace of a positive operator taken as its
diagonal sum in a Hilbert basis; the facts (T1)--(T4) themselves,
\Cref{rem:duality_confinement:analysis-operator} and its use in
\Cref{sec:domination} are not formally verified, and
\Cref{app:formalization} lists the correspondence item by item.

\paragraph{Possible instances are open problems.}
The instances discussed in \Cref{sec:discussion} indicate where the
framework might apply.  None is claimed as a result: in each case the
domination records would have to be produced by an argument specific
to the instance.

\paragraph{No novelty is claimed for classical operator theory.}
The Loewner order \citep{Loewner1934,Bhatia1997,HornJohnson2013},
trace ideals \citep{Simon2005,ReedSimon1980}, Markov's inequality
\citep{Folland1999}, the arithmetic--geometric mean inequality
\citep{HardyLittlewoodPolya1952} and Douglas's factorization theorem
\citep{Douglas1966} are classical, and the proofs given here are
standard.  The contribution of the paper is the formulation: the
involutive object ledger, the odd readout and its Gram operator
\(\AX\), the bookkeeping of domination records with explicit defects,
and the reduction of fixed-locus confinement to trace estimates.
 \section{Discussion}
\label{sec:discussion}

\paragraph{A worked instance: self-adjointness.}
The definitions have a transparent instance in matrix analysis.  Let
\(X\) be the space \(M_d(\mathbb C)\) of complex \(d\times d\) matrices,
\(\mu\) a finite positive Borel measure on \(X\) with
\(\int_X\|T\|_{\mathrm{HS}}^2\,d\mu(T)<\infty\), and \(J(T)=T^*\).
Regard \(Y=M_d(\mathbb C)\) as a real Hilbert space with the inner
product \(\langle A,B\rangle=\operatorname{Re}\operatorname{tr}(A^*B)\),
let \(\Jiso(A)=A^*\), which is a real-linear isometric involution, and
let \(\psi\) be the identity.  Then
\[
  \psim(T)=\tfrac12\,(T-T^*)
\]
is the skew-Hermitian part of \(T\), the fixed locus is the set of
Hermitian matrices, and \(\psim\) is qualitatively separating.  Since
the Hermitian part \(\tfrac12(T+T^*)\) is the Hermitian matrix nearest
to \(T\) in the Hilbert--Schmidt norm, \(\dist(T,\Fix(J))=\|\psim(T)\|\)
and \(\psim\) is quantitatively separating with \(m(\varepsilon)=\varepsilon\).
\Cref{thm:duality_confinement:master-theorem} then says: \(\mu\) is
carried by the Hermitian matrices if and only if the Gram operator of
the skew-Hermitian parts is dominated by positive operators of
vanishing trace.  \Cref{thm:duality_confinement:quantitative-confinement}
bounds the \(\mu\)-mass of matrices at distance at least \(\varepsilon\)
from the Hermitian ones by \(\varepsilon^{-2}\trace\AX\).  Since
\(\trace\AX=\int_X\|\psim\|^2\,d\mu\) can be computed directly here,
the instance illustrates the definitions but gains nothing from the
domination formulation; a substantive use would need a situation in
which \(\AX\) is known only through a bridge to independently computed
operators.

\paragraph{Curves over finite fields.}
For a smooth projective curve over a finite field, the analogue of the
Riemann hypothesis is a theorem of Weil \citep{Weil1948a,Weil1948b}.
Stepanov introduced an elementary polynomial method and used it to
bound the number of points on hyperelliptic curves over prime fields
\citep{Stepanov1969}; Bombieri, building on that method, gave an
elementary proof of Weil's theorem for all curves \citep{Bombieri1974}.
The zeta function of a curve of genus \(g\) over \(\mathbb F_q\) has
\(2g\) inverse roots \(\alpha\), permuted by the involution
\(\alpha\mapsto q/\overline{\alpha}\) coming from the functional
equation; its fixed locus is the circle \(|\alpha|=\sqrt q\).  The
corresponding ledger is therefore a finite weighted one.  Since the
conclusion is known, trace-vanishing records exist trivially (the zero
sequence), so the meaningful question is a different one: whether the
Stepanov--Bombieri counting argument can be recast as a bridge
\(\AX\preceq\Kminus+E\) in which \(\Kminus\) and \(E\) are computed
from point counts rather than from the roots.  We have not carried
this out.

\paragraph{Relation to classical operator theory.}
Stripped of the framework, \Cref{thm:duality_confinement:master-theorem}
is the familiar fact that a positive trace-class operator dominated by
operators of arbitrarily small trace is zero, applied to the Gram
operator of an odd observable.  What the formulation adds is a fixed
interface: an application must specify the involution, the equivariant
readout and the domination records, and the conclusion is then read
off without further analysis.

\paragraph{Further directions.}
The natural direction is to find instances in which domination
records can be produced by estimates, for instance through the
factorization criterion of \Cref{sec:domination}.
 \section{Conclusion}
\label{sec:conclusion}

For an involution \(J\) of a measure space, a linear isometric
involution of a real Hilbert space, and a strongly measurable
equivariant readout with square-integrable odd part \(\psim\), the Gram operator
\(\AX=\int_X\psim\psim^*\,d\mu\) is positive and trace-class with
\(\trace\AX=\int_X\|\psim\|^2\,d\mu\).  Domination
\(\AX\preceq B_n\) by positive trace-class operators with
\(\trace B_n\to0\) forces \(\AX=0\), hence \(\psim=0\) almost
everywhere, hence---when \(\psim\) is separating---confinement of the
mass to the fixed locus \(\Fix(J)\); conversely, confinement makes the
zero sequence a valid set of records.  The same conclusion follows from
domination on finite windows up to positive tails when the transported
window bounds and the tails have traces tending to zero, and from
two-defect budgets for \(\AX\) whose optimized trace
\((\sqrt a+\sqrt b)^2\) tends to zero.  When the odd readout is
quantitatively separating with a positive modulus \(m\) on a metric
space, the outer measure of the points at distance at least
\(\varepsilon\) from \(\Fix(J)\) is at most
\(\trace\AX/m(\varepsilon)^2\).  Douglas factorization identifies exact
domination with contractive factorization of the analysis operator of
\(\psim\).

In the Six Birds framework these results are the mathematical side of
the structural law \(\GamSDTC\), which asserts that closures with
genuine involutive self-duality supply such domination records.  The
law is a hypothesis; the theorems of this paper say exactly what
follows from it and what it amounts to.  The results, for real Hilbert spaces of arbitrary dimension, and the
abstract squeeze are formally verified in Lean~4 with Mathlib
(\Cref{app:formalization}).

\appendix
\crefalias{section}{appendix}
\section{Formal verification}
\label{app:formalization}

Every theorem, lemma and proposition of this paper, and the abstract
squeeze of \Cref{rem:duality_confinement:abstract-cone}, has been
formally verified in the Lean~4 proof assistant (version 4.28.0) using
the Mathlib library (revision
\nolinkurl{8f9d9cff6bd728b17a24e163c9402775d9e6a365}), for real Hilbert
spaces of arbitrary dimension.  This appendix states how the objects of
the paper are represented, what is checked only in the text, and which
Lean theorem proves each numbered result.

\subsection*{Representation}

The response space \(Y\), the spaces \(H,K_1,K_2\) of
\Cref{thm:duality_confinement:douglas-domination} and the window
spaces \(Y_n\) of \Cref{thm:duality_confinement:exhaustive-squeeze}
are complete real inner product spaces; no finite dimension or
separability is assumed.  The operator \(\Jiso\) is a linear isometric
equivalence of \(Y\) that is an involution, the readout \(\psi\) is
\(\mu\)-almost everywhere strongly measurable, and the involution \(J\)
need not be measurable.  The ledger \(\AX\) is the Bochner integral of
\(x\mapsto\psim(x)\psim(x)^*\) in the operator norm.  The trace of a
positive operator \(T\) is the extended sum
\(\sum_i\langle Te_i,e_i\rangle\in[0,\infty]\) over a Hilbert basis
\((e_i)\); it is proved independent of the basis, and a positive
trace-class operator is represented as a positive bounded operator whose
trace is finite.  The order \(A\preceq B\) is the inequality
\(\langle Ah,h\rangle\leq\langle Bh,h\rangle\) of quadratic forms, which
for the self-adjoint operators of the paper is the Loewner order.  In
\Cref{thm:duality_confinement:quantitative-confinement} the distance to
\(\Fix(J)\) is the extended distance, which is \(+\infty\) when
\(\Fix(J)\) is empty.

The formal proofs do not use the trace-class facts (T1)--(T4) of
\Cref{sec:anti_invariant_ledger}; they work directly with the operator
norm Bochner integral and the extended trace.

\subsection*{What is checked only in the text}

The following are not formally verified: that \(x\mapsto\psim(x)\psim(x)^*\)
is Bochner integrable in the trace class \(\mathcal S_1(Y)\), and the
facts (T1)--(T4) themselves, including the identification, for positive
operators, of finite extended trace with membership in the trace class
and of the extended sum with the trace (standard facts recalled in
\Cref{sec:anti_invariant_ledger}); the comparison of quantitative and
qualitative separation after
\Cref{def:duality_confinement:separating-readout};
\Cref{rem:duality_confinement:analysis-operator} and the discussion of
exact bridges as factorizations in \Cref{sec:domination}; the
Hilbert--Schmidt reading of
\Cref{prop:duality_confinement:optimized-trace-budget}; the worked
examples; and the framework material of \Cref{sec:framework}.  The
structural law \(\GamSDTC\) is a hypothesis and is not formalized as a
theorem.

\Needspace{16\baselineskip}
\subsection*{Correspondence}

{\small
\begin{longtable}{@{}>{\raggedright\arraybackslash}p{0.2\linewidth}
  >{\raggedright\arraybackslash}p{0.75\linewidth}@{}}
\caption{Numbered results and the Lean theorems that prove them for
real Hilbert spaces of arbitrary dimension; names are relative to the
namespace \texttt{SixBirdsDualityConfinement.DualityConfinement}.
\Cref{def:duality_confinement:separating-readout,def:duality_confinement:completed-domination-bridge,def:duality_confinement:exhaustive-moving-ledger,def:duality_confinement:defected-budget}
record hypotheses; in Lean they appear as the hypotheses of the
theorems listed for the results that use them.  The modules
\texttt{ConcreteFinite} and \texttt{ConcreteMeasured} contain
separate proofs of the finite-dimensional cases with the
linear-algebra trace.}\label{tab:dc-formalization-coverage}\\
\toprule
Result & Lean theorems \\
\midrule
\endfirsthead
\toprule
Result & Lean theorems (continued) \\
\midrule
\endhead
\bottomrule
\endfoot
\Cref{def:duality_confinement:involutive-object-ledger,def:duality_confinement:anti-invariant-ledger}
  & \leanid{ConcreteCore.pminus}, \leanid{ConcreteCore.psiMinus}, \leanid{ConcreteHilbert.ledger}, \leanid{ConcreteHilbert.PositiveTraceClass}, \leanid{ConcreteHilbert.positiveTraceClass_basis_iff}, \leanid{ConcreteHilbert.psiMinus_aestronglyMeasurable}, \leanid{ConcreteCore.psiMinus_equivariant}, \leanid{ConcreteCore.psiMinus_fixed} \\ \addlinespace
\Cref{lem:duality_confinement:trace-identity}
  & \leanid{ConcreteHilbert.ledger_positive}, \leanid{ConcreteHilbert.ledger_inner}, \leanid{ConcreteHilbert.ledger_trace}, \leanid{ConcreteHilbert.ledger_traceClass} \\ \addlinespace
\Cref{thm:duality_confinement:separation-confinement}
  & \leanid{ConcreteHilbert.ledger_zero_iff}, \leanid{ConcreteHilbert.separation_confinement}, \leanid{ConcreteHilbert.atom_readout_of_zero}, \leanid{ConcreteHilbert.atom_fixed_of_zero} \\ \addlinespace
\Cref{thm:duality_confinement:quantitative-confinement}
  & \leanid{ConcreteHilbert.quantitative_threshold} (arbitrary subsets, outer measure), \leanid{ConcreteHilbert.quantitative_distance} \\ \addlinespace
\Cref{thm:duality_confinement:douglas-domination}
  & \leanid{ConcreteHilbert.douglas_factorization} \\ \addlinespace
\Cref{thm:duality_confinement:master-theorem}
  & \leanid{ConcreteHilbert.ledger_master}, \leanid{ConcreteHilbert.ledger_master_fixed}; converse \leanid{ConcreteHilbert.ledger_records_iff}, \leanid{ConcreteHilbert.ledger_zero_of_ae}, \leanid{ConcreteHilbert.ledger_zero_of_ae_fixed}, \leanid{ConcreteHilbert.zero_budget_records} \\ \addlinespace
\Cref{rem:duality_confinement:abstract-cone}
  & \leanid{AbstractSqueeze.abstract_trace_squeeze} \\ \addlinespace
\Cref{thm:duality_confinement:exhaustive-squeeze}
  & \leanid{ConcreteHilbert.ledger_exhaustive}, \leanid{ConcreteHilbert.ledger_exhaustive_fixed} \\ \addlinespace
\Cref{prop:duality_confinement:optimized-trace-budget}
  & \leanid{ConcreteHilbert.scalar_budget_infimum}, \leanid{ConcreteHilbert.optimized_budget_infimum}, \leanid{ConcreteHilbert.ledger_optimized_budget}, \leanid{ConcreteHilbert.ledger_budget_sequence}, \leanid{ConcreteHilbert.ledger_budget_sequence_fixed}; attainment at \(t=\sqrt{b/a}\) from \leanid{youngTrace_optimum} in \texttt{SixBirdsNeedles.MainCore} \\
\end{longtable}
}

The general Hilbert-space theorems are derived from the operator
ledger, positive trace, confinement, exhaustion, trace budget and
Douglas factorization results of the companion library
\texttt{SixBirdsNeedles.MainCore} \citep{TsiokosNeedleKiller2026}.

\subsection*{Trust base and availability}

Every theorem listed above depends only on the three standard axioms
of Lean's type theory with classical logic (\texttt{propext},
\texttt{Classical.choice}, \texttt{Quot.sound}); no \texttt{sorry} and
no additional axioms occur.  The development is available at
\url{https://github.com/ioannist/six-birds-duality-confinement}.  The
same repository contains an earlier formulation of these results over
an abstract ``typed positive cone'', in which the trace identity,
Douglas factorization and the arithmetic--geometric mean step are
supplied as hypotheses; the companion Riemann-hypothesis development
builds on that formulation.  The claims of the present paper rest only
on the theorems listed in \Cref{tab:dc-formalization-coverage}.
 \section{Version history}
\label{app:version-history}

Every numbered result of Version~1 keeps its number and place in later
versions, and its name except where noted below.

\subsection*{Version 3 (October 1, 2026)}

\begin{itemize}
  \item The results are now stated and proved for positive trace-class
    operators on an arbitrary real Hilbert space.  Version~2 stated
    them for an abstract ``typed positive cone'' and described the
    trace-class case as an open extension.  The trace identity
    (\Cref{lem:duality_confinement:trace-identity}) is now proved from
    Bochner integration in the trace class; Version~2 took it as part
    of the data.
  \item The readout in \Cref{def:duality_confinement:involutive-object-ledger}
    is required to be strongly measurable; together with
    \(\int_X\|\psim\|^2\,d\mu<\infty\) this makes the Bochner integral
    defining \(\AX\) exist.
    \Cref{def:duality_confinement:anti-invariant-ledger,def:duality_confinement:completed-domination-bridge,def:duality_confinement:exhaustive-moving-ledger,def:duality_confinement:defected-budget}
    now specify the operators involved (positive trace-class operators,
    bounded transports between Hilbert spaces).
  \item The confinement conclusions are stated as ``\(Jx=x\) for
    \(\mu\)-almost every \(x\)'', which needs no measurability of
    \(\Fix(J)\); Version~2 assumed \(X\setminus\Fix(J)\) measurable.
    \Cref{thm:duality_confinement:quantitative-confinement} is stated
    for outer measure, so the far set need not be measurable, and it
    holds for any displacement function \(\delta\) in place of the
    distance, provided \(\delta(x)\geq\varepsilon\) implies
    \(\|\psim(x)\|\geq m(\varepsilon)>0\).
  \item \Cref{thm:duality_confinement:master-theorem} now includes its
    converse: trace-vanishing domination records exist if and only if
    the odd readout vanishes almost everywhere.  The abstract,
    \Cref{sec:intro,sec:framework,sec:master_theorem,sec:scope} now
    state the consequences: the theorem does not make confinement
    easier to prove, its role is to fix the form of a bridge argument,
    and SDTC cannot hold for all involutive ledgers (the
    complex-conjugation example is a counterexample).  Version~2 did
    not state the converse or these consequences.
  \item \Cref{thm:duality_confinement:douglas-domination} is stated for
    bounded operators between real Hilbert spaces, with a complete proof
    and a precise account of Douglas's original formulation; its name
    changed from ``Douglas factorization, typed-cone form'' to
    ``Douglas factorization''.  The new
    \Cref{rem:duality_confinement:analysis-operator} identifies \(\AX\)
    with \(V^*V\) for the analysis operator \(V\) of the odd readout.
  \item \Cref{prop:duality_confinement:optimized-trace-budget} is proved
    over the real numbers, including the cases \(a=0\) or \(b=0\), and
    now includes the consequence for sequences of budgets.  Version~2
    treated the arithmetic--geometric mean step as a hypothesis on an
    abstract scalar type.
  \item The formal verification was extended: every theorem, lemma
    and proposition is now proved in Lean~4 with Mathlib for real Hilbert
    spaces of arbitrary dimension, as is the abstract squeeze
    (\Cref{rem:duality_confinement:abstract-cone}).
    \Cref{app:formalization} was rewritten accordingly.
  \item The exposition was rewritten for readers outside the framework:
    the introduction states the results with their hypotheses,
    \Cref{sec:framework} explains the framework terms, and worked
    examples were added (complex conjugation, self-adjointness of
    matrices).  The Version~2 reference to a ``V-Differential''
    classification, which is not defined in the cited literature, was
    removed, as was a paragraph on charge conjugation, CPT and gauge
    symmetries that made no claim.  The discussion of curves over finite
    fields now states
    that the Riemann hypothesis there is a theorem of Weil, attributes
    the elementary method to Stepanov and its general form to Bombieri,
    and presents that setting as a test case rather than a prediction.
  \item Cross-references now name the kind of result they point to;
    in Version~2 references to definitions, lemmas, propositions and
    remarks were printed as ``Theorem''.  The DOI of
    \citet{Stepanov1969} in the bibliography was corrected.
\end{itemize}

\subsection*{Version 2 (September 28, 2026)}

\begin{itemize}
  \item The hypotheses of the measure-theoretic conclusions were made
    explicit: integrability of \(\|\psim\|^2\), the trace identity, and
    the passage from almost-everywhere vanishing to a measure
    statement.
  \item The defected obstruction budget was restated as a family of
    domination records for the same operator \(\AX\), one for each
    \(t>0\), which is what the trace bound of
    \Cref{prop:duality_confinement:optimized-trace-budget} requires.
  \item The description of the formal verification was corrected to
    distinguish the results derived in Lean from those whose inputs are
    supplied as hypotheses, and a numerical form of the trace squeeze
    was added.
  \item The description of the companion Riemann-hypothesis paper was
    corrected to state which of its hypotheses are supplied.
\end{itemize}

\clearpage
\bibliographystyle{plainnat}
\begin{thebibliography}{18}
\providecommand{\natexlab}[1]{#1}
\providecommand{\url}[1]{\texttt{#1}}
\expandafter\ifx\csname urlstyle\endcsname\relax
  \providecommand{\doi}[1]{doi: #1}\else
  \providecommand{\doi}{doi: \begingroup \urlstyle{rm}\Url}\fi

\bibitem[Bhatia(1997)]{Bhatia1997}
Rajendra Bhatia.
\newblock \emph{Matrix Analysis}, volume 169 of \emph{Graduate Texts in Mathematics}.
\newblock Springer, New York, 1997.
\newblock ISBN 978-0-387-94846-1.
\newblock \doi{10.1007/978-1-4612-0653-8}.

\bibitem[Bombieri(1974)]{Bombieri1974}
Enrico Bombieri.
\newblock Counting points on curves over finite fields (d'apr{\`e}s {S}. {A}. {Stepanov}).
\newblock In \emph{S{\'e}minaire Bourbaki, 25e ann{\'e}e (1972/1973), Exp. no. 430}, volume 383 of \emph{Lecture Notes in Mathematics}, pages 234--241. Springer-Verlag, Berlin, 1974.

\bibitem[Douglas(1966)]{Douglas1966}
R.~G. Douglas.
\newblock {On Majorization, Factorization, and Range Inclusion of Operators on {Hilbert} Space}.
\newblock \emph{Proceedings of the American Mathematical Society}, 17\penalty0 (2):\penalty0 413--415, 1966.
\newblock \doi{10.1090/S0002-9939-1966-0203464-1}.

\bibitem[Folland(1999)]{Folland1999}
Gerald~B. Folland.
\newblock \emph{Real Analysis: Modern Techniques and Their Applications}.
\newblock Wiley, New York, second edition, 1999.
\newblock ISBN 978-0-471-31716-6.

\bibitem[Hardy et~al.(1952)Hardy, Littlewood, and P{\'o}lya]{HardyLittlewoodPolya1952}
G.~H. Hardy, J.~E. Littlewood, and G.~P{\'o}lya.
\newblock \emph{Inequalities}.
\newblock Cambridge University Press, Cambridge, second edition, 1952.
\newblock ISBN 978-0-521-05206-1.

\bibitem[Horn and Johnson(2013)]{HornJohnson2013}
Roger~A. Horn and Charles~R. Johnson.
\newblock \emph{Matrix Analysis}.
\newblock Cambridge University Press, Cambridge, second edition, 2013.
\newblock ISBN 978-0-521-54823-6.

\bibitem[L{\"o}wner(1934)]{Loewner1934}
Karl L{\"o}wner.
\newblock {\"U}ber monotone {M}atrixfunktionen.
\newblock \emph{Mathematische Zeitschrift}, 38\penalty0 (1):\penalty0 177--216, 1934.
\newblock \doi{10.1007/BF01170633}.

\bibitem[Reed and Simon(1980)]{ReedSimon1980}
Michael Reed and Barry Simon.
\newblock \emph{Methods of Modern Mathematical Physics {I}: Functional Analysis}.
\newblock Academic Press, New York, revised and enlarged edition, 1980.
\newblock ISBN 978-0-12-585050-6.

\bibitem[Simon(2005)]{Simon2005}
Barry Simon.
\newblock \emph{Trace Ideals and Their Applications}, volume 120 of \emph{Mathematical Surveys and Monographs}.
\newblock American Mathematical Society, Providence, RI, second edition, 2005.
\newblock ISBN 978-0-8218-3581-4.
\newblock \doi{10.1090/surv/120}.

\bibitem[Stepanov(1969)]{Stepanov1969}
S.~A. Stepanov.
\newblock On the number of points of a hyperelliptic curve over a finite prime field.
\newblock \emph{Mathematics of the USSR-Izvestiya}, 3\penalty0 (5):\penalty0 1103--1114, 1969.
\newblock \doi{10.1070/IM1969v003n05ABEH000834}.

\bibitem[Titchmarsh(1986)]{TitchmarshHeathBrown1986}
E.~C. Titchmarsh.
\newblock \emph{The Theory of the {Riemann} Zeta-Function}.
\newblock Oxford University Press, Oxford, second edition, 1986.
\newblock ISBN 978-0-19-853369-6.
\newblock Revised by D. R. Heath-Brown.

\bibitem[Tsiokos(2026{\natexlab{a}})]{TsiokosFoundationsI2026}
Ioannis Tsiokos.
\newblock {Six Birds: Foundations of Emergence Calculus}.
\newblock \href{https://arxiv.org/abs/2602.00134}{arXiv:2602.00134}, 2026{\natexlab{a}}.

\bibitem[Tsiokos(2026{\natexlab{b}})]{TsiokosFoundationsII2026}
Ioannis Tsiokos.
\newblock {Six Birds Foundations II: Admissibility, Meta-Theory, and the Exact-Six Program}.
\newblock \doi{10.5281/zenodo.19672278}, 2026{\natexlab{b}}.

\bibitem[Tsiokos(2026{\natexlab{c}})]{TsiokosFoundationsIII2026}
Ioannis Tsiokos.
\newblock {Six Birds Foundations III: A Finite Audited Interaction Calculus for SBT}.
\newblock \doi{10.5281/zenodo.20124595}, 2026{\natexlab{c}}.

\bibitem[Tsiokos(2026{\natexlab{d}})]{TsiokosNeedleKiller2026}
Ioannis Tsiokos.
\newblock Emergence {IS} the needle killer.
\newblock \doi{10.5281/zenodo.23086988}, 2026{\natexlab{d}}.
\newblock Version 4, October 1, 2026.

\bibitem[Tsiokos(2026{\natexlab{e}})]{TsiokosRHviaSDTC2026}
Ioannis Tsiokos.
\newblock {Riemann Hypothesis} via self-dual trace confinement: A conditional closure.
\newblock \doi{10.5281/zenodo.23086975}, 2026{\natexlab{e}}.
\newblock Version 3, October 1, 2026.

\bibitem[Weil(1948{\natexlab{a}})]{Weil1948a}
Andr{\'e} Weil.
\newblock \emph{Sur les courbes alg{\'e}briques et les vari{\'e}t{\'e}s qui s'en d{\'e}duisent}, volume 1041 of \emph{Actualit{\'e}s scientifiques et industrielles}.
\newblock Hermann, Paris, 1948{\natexlab{a}}.
\newblock Publications de l'Institut de math{\'e}matique de l'Universit{\'e} de Strasbourg, no. 7.

\bibitem[Weil(1948{\natexlab{b}})]{Weil1948b}
Andr{\'e} Weil.
\newblock \emph{Vari{\'e}t{\'e}s ab{\'e}liennes et courbes alg{\'e}briques}, volume 1064 of \emph{Actualit{\'e}s scientifiques et industrielles}.
\newblock Hermann, Paris, 1948{\natexlab{b}}.
\newblock Publications de l'Institut de math{\'e}matique de l'Universit{\'e} de Strasbourg, no. 8.

\end{thebibliography}


\end{document}
